Após configurar e instalar todas as ferramentas necessárias, é comum começar 
com um simples ``Olá, mundo!''.
\\\\
O programa que iremos fazer exibe o texto ``Olá, mundo!'' na tela e embora seja simples, é 
um ponto de partida essencial para começar a entender não somente como a linguagem C 
funciona, mas qualquer outra linguagem de programação.
\\\\
Para criar o seu primeiro programa, siga estes passos:

\begin{enumerate}
    \item Crie um diretório para guardar o seu projeto.
    \item Crie um arquivo chamado ``main.c'' onde será escrito o código.
    \item Abra o arquivo ``main.c'' no seu editor de código e escreva o seguinte:

        \begin{verbatim}
        #include <stdio.h>

        int main(void) {
            printf("Olá, mundo!\n");
            return 0;
        }
        \end{verbatim}
    \item Salve o arquivo de texto e compile o código com o GCC usando o seguinte comando no seu terminal:
        \begin{verbatim}
        gcc main.c
        \end{verbatim}
    \item Agora, execute o programa compilado com o seguinte comando:
        \begin{verbatim}
        # Para sistemas Linux
        ./a.out
        # Para sistemas Windows
        ./a.exe
        \end{verbatim}
\end{enumerate}

Após executar o comando acima, você já deve ver o texto ``Olá, mundo!'' na sua 
tela. Mas antes de prosseguir, você provavelmente está se perguntando o que 
exatamente cada uma das linhas de código que você escreveu fazem e como 
exatamente elas fazem esse texto aparecer na sua tela. Por isso, vou explicar passo a 
passo o que significa cada linha do código acima.

\begin{verbatim}
    #include <stdio.h>
\end{verbatim}

Na linha número 1, temos o ``include'', uma diretiva que inclui bibliotecas 
externas no nosso programa. Aqui, estamos incluindo a biblioteca ``stdio.h''. 
O conceito de bibliotecas ainda vai ser melhor explicado futuramente, mas você pode 
entender como se cada biblioteca fosse um tipo de extensão com funções extras.

\begin{verbatim}
    int main(void) {
\end{verbatim}

Na linha número 2, temos o início de uma função. Nessa linha, há o uso 
de quatro palavras/caracteres importantes:

\begin{itemize}
    \item \textbf{"int":} Define o que a função deve retornar ao ser finalizada, 
        neste caso, um número inteiro.
    \item \textbf{"main":} É o nome que foi escolhido para o bloco de código 
        principal. Apesar de não ser necessariamente obrigatório usar o nome 
        \textbf{main} para a sua função principal, acabar por ser uma boa prática 
        a ser seguida.
    \item \textbf{"(void)":} Define quais parâmetros/argumentos são necessários para 
        chamar a função. Neste caso, a palavra \textbf{void} indica que a função não 
        precisa de nenhum valor para ser chamada.
    \item \textbf{"\{":} Marca o início de um bloco de código de uma função.
\end{itemize}

\begin{verbatim}
    printf("Olá, mundo!\n");
\end{verbatim}

Na linha número 3, encontramos a diretiva ``printf'', responsável por exibir 
informações no terminal. Esta diretiva na verdade não existe como uma funcionalidade 
padrão na linguagem C, e está sendo implementada através da biblioteca 
``stdio.h'', mencionada na linha 1.
\\\\
Neste caso a informação que está sendo passada para a função é o texto "Olá, 
mundo!\textbackslash n" que sempre deve ser colocado entre os parênteses e as aspas, 
você pode ter percebido também a presença do caracter ``\textbackslash n'' que não 
aparece quando a mensagem é exibida. Esse caracter é um caracter especial e tem 
como função fazer com que naquele espaço da mensagem o terminal desça uma 
linha abaixo antes de continuar.

\begin{verbatim}
    return 0;
\end{verbatim}

Na linha número 4, encontramos o comando ``return''. Em uma função, essa diretiva
indica o valor que a função deve retornar quando terminar sua execução e normalmente, 
mas nem sempre, é acompanhada de um valor que está escrito como um sufixo da diretiva. 
\\\\
No caso da função ``main'', que é a função principal em um programa C, o valor 0 
geralmente é usado como uma boa prática para indicar que o programa foi finalizado sem 
erros.

\begin{verbatim}
    }
\end{verbatim}

A linha número 5 é a chave de fechamento ``\}'', que indica o fim do bloco de 
código da função ``main''. Tudo o que está contido entre as chaves forma o corpo 
da função, o compilador (GCC) utiliza dessas chaves de abertura e fechamento como uma 
forma de entender quando um código é pertencente a uma função.
\\\\
Caso tenha parecido um pouco confuso a primeira vista, acredito que com um pouco de 
prática você vai começar a escrever esse tipo de código sem sequer pensar direito 
sobre isso. E no geral, essa estrutura tende a se repetir bastante, então caso não 
tenha entendido algo muito bem vale a pena voltar e reler.

\subsubsection{Exercícios}

\textbf{É de extrema importância que você realize os exercícios aqui elaborados para fixar o 
conteúdo.}

\begin{enumerate}
    \item{Modifique a mensagem exibida na tela pela função ``printf''.}
    \item{Tente alterar a posição do caracter especial ``\textbackslash n'' e veja o que acontece.}
    \item{Mudar o nome da função para qualquer outra coisa e veja qual o erro retornado.}
    \item{Remova a linha de código que inclui a biblioteca ``stdio.h'' e veja qual o erro retornado.}
\end{enumerate}
